\clearpage
\SourceLeaf{454}
\begin{center}\large [ALLGEMEINSTE LÖSUNG\\DES PROBLEMS DER ABWICKLUNG VON FLÄCHEN.]\end{center}
\noindent\hspace*{1.2em}Die ganz allgemeine Lösung des Problems der Abwicklung von Flächen lautet:

1) Die Natur einer krummen Fläche wird dadurch bestimmt, dass $z$ als endliche Funktion von $x,y$ dargestellt wird. Die Differentiation ergebe
\[
dz=pdx+qdy.
\]

2) Man setze
\[
t=\frac{1}{\BareRoot{(pp+qq+1)}}
\]
\noindent und nehme an, dass die Differentiation einer beliebigen Funktion von $t,T$ den Ausdruck
\[
udt+UdT,
\]
\noindent ergibt, wobei $u,U$ gegebene Funktionen von $t$ und $T$ sein werden.

Durch Elimination werden also $T,U$ zu Funktionen von $t,u$ und, wenn
\[
\operatorname{tang}u=\frac{q}{p}
\]
\noindent gesetzt wird, zu Funktionen von $x$ und $y$.

3) Man setze nun
\[
\frac{Q}{P}=\operatorname{tang}U,\qquad
T\BareRoot{(1+PP+QQ)}=1,
\]
\noindent so werden $P,Q$ zu Funktionen von $x$ und $y$.

